\clearpage
\SourceLeaf{149}
\begin{center}
[IV.]

{\large [KRITISCHE BEMERKUNGEN ZUR METHODE DER KLEINSTEN QUADRATE.]}
\end{center}



\Needspace{6\baselineskip}
\begin{center}
\textsc{Gauss} an \textsc{Olbers}. Göttingen, 22. Februar 1819.
\end{center}

\noindent . . . . . . . . . . . Auch bin ich jetzt mit einer neuen Begründung der sogenannten Methode der kleinsten Quadrate beschäftigt. Meine erste Begründung setzt voraus, dass die Wahrscheinlichkeit des Beobachtungsfehlers $x$ durch $e^{-hhxx}$ dargestellt werde, wo denn jene Methode nach aller Strenge und in allen Fällen die wahrscheinlichsten Resultate gibt. Ist das Gesetz der Fehler unbekannt, so ist es \textls{unmöglich} die wahrscheinlichsten Resultate aus \textls{schon gemachten} Beobachtungen anzugeben. \textsc{Laplace} hat die Sache von einer verschiedenen Seite angesehen und ein Princip gewählt, welches auch auf die M.\,d.\,kl.\,Q.\ führt\HistoricalNote{*)}{*)}{und zwar unabhängig von dem Fehlergesetz.}, wenn die Anzahl der Beobachtungen unendlich gross ist. Allein bei einer mässigen Anzahl Beobachtungen bleibt man, wenn das Fehlergesetz unbekannt, ganz im Dunkeln, und \textsc{Laplace} weiss auch selbst für diesen Fall nichts besseres zu sagen, als dass man die Meth.\ der kl.\ Quadrate auch hier anwenden möge, weil sie bequeme Rechnung gewähre. Ich habe jetzt gefunden, dass bei der Wahl eines etwas andern Princips als das \textsc{Laplace}sche (und zwar eines solchen, wo \textls{niemand} in Abrede stellen kann, dass man wenigstens eben so gut zu dessen Annahme befugt sei als zu dem von \textsc{Laplace}, und welches, meiner Meinung nach, \textls{jeder nicht im Voraus Eingenommene} für natürlicher erklären muss als das \textsc{Laplace}sche) man alle jene Vor%
